% ==========================================
% BAB V IMPLEMENTASI
% ==========================================
\cleardoublepage
\chapter{IMPLEMENTASI}
\label{chap:implementasi}

\section{Struktur Kode dan Lingkungan}
\label{sec:struktur-kode}
Seluruh kode berada dalam satu paket bernama \texttt{gfd}.
Paket ini dipasang sekali dengan \texttt{pip install -e code/}, lalu tiap bagiannya dipanggil sebagai modul.
Susunan paket dapat dilihat pada Listing~\ref{lst:struktur-gfd}.

\begin{minipage}{\textwidth}
\begin{lstlisting}[caption={Susunan paket gfd}, label={lst:struktur-gfd}, frame=single, captionpos=t, language=]
gfd/
    config.py
    dataset/
    selection/
    models/
    training/
    evaluation/
    experiments/
\end{lstlisting}
\end{minipage}

Berkas \texttt{config.py} memegang hal yang harus disepakati semua subpaket.
Subpaket \texttt{dataset/} mengambil data dari tiap sumber dan menggabungkannya menjadi tabel per jam.
Subpaket \texttt{selection/} memilih fitur yang dipakai model.
Subpaket \texttt{models/} hanya memuat definisi estimator, sedangkan pelatihannya ada di \texttt{training/}.
Subpaket \texttt{evaluation/} menghitung metrik dan menyusun tabel hasil, dan \texttt{experiments/} menjalankan matriks eksperimen.
Pembagian ini memisahkan bagian yang berubah karena sebab berbeda: definisi model, prosedur pelatihan, dan tabel hasil.

Lingkungan yang dipakai adalah Python 3.13.7, dengan batas bawah 3.12 karena numpy 2.5.2 tidak dapat dipasang di bawahnya.
Seluruh versi pustaka dipatok tepat pada satu berkas \texttt{requirements.txt}.
Versi yang dipatok adalah versi yang menghasilkan tabel dan hasil pada tugas akhir ini, jadi pemutakhiran pustaka tidak dilakukan tanpa alasan yang dicatat.
Daftar lengkapnya dapat dilihat pada Tabel~\ref{tab:pustaka}.

\begin{table}[!t]
  \caption{Pustaka Python yang dipakai beserta versinya}
  \label{tab:pustaka}
  \small
  \begin{tabular}{ | l | l | p{0.42\textwidth} | }
    \hline
    Pustaka & Versi & Kegunaan \\
    \hline
    \multicolumn{3}{|l|}{\textit{Pengolahan data}} \\
    \hline
    pandas & 3.0.5 & Tabel data \\
    numpy & 2.5.2 & Operasi larik numerik \\
    scipy & 1.17.1 & Fungsi statistik \\
    pyarrow & 25.0.1 & Membaca dan menulis berkas \textit{parquet} \\
    scikit-learn & 1.8.0 & Model klasik pembanding, penskalaan, dan metrik \\
    matplotlib & 3.10.8 & Gambar \\
    \hline
    \multicolumn{3}{|l|}{\textit{Akuisisi dan pembacaan berkas sumber}} \\
    \hline
    requests & 2.34.2 & Permintaan HTTP ke arsip MERLIN dan NASA POWER \\
    truststore & 0.10.4 & Sertifikat sistem, tanpa ini unduhan MERLIN gagal \\
    cdsapi & 0.7.7 & Permintaan ke \textit{Climate Data Store} untuk ERA5 \\
    xarray & 2026.7.0 & Membaca berkas keluaran ERA5 \\
    netCDF4 & 1.7.4 & Pembaca berkas NetCDF \\
    openpyxl & 3.1.5 & Membaca berkas Excel data PLN \\
    \hline
    \multicolumn{3}{|l|}{\textit{Kuantum dan optimasi}} \\
    \hline
    qiskit & 2.5.2 & Pembangunan sirkuit kuantum \\
    qiskit-machine-learning & 0.9.1 & Komponen pembelajaran mesin kuantum \\
    qiskit-aer & 0.17.2 & Simulator sirkuit Qiskit \\
    pennylane & 0.45.1 & Kerangka kerja QNN \\
    pennylane-lightning & 0.45.0 & Simulator cepat untuk PennyLane \\
    torch & 2.14.0 & Perhitungan gradien untuk PennyLane \\
    \hline
  \end{tabular}
\end{table}

Terdapat dua pustaka kuantum, dan ini disengaja.
Qiskit hanya menyediakan \mbox{\textit{parameter-shift}} untuk perhitungan gradien, dan cara itu terlalu lambat untuk pelatihan.
PennyLane dipakai untuk pelatihan, sedangkan Qiskit dipertahankan sebagai pembanding untuk memeriksa kedua sirkuit menghasilkan hal yang sama.
Keduanya sudah dipatok versinya, tetapi belum terpakai.
Dari tujuh nama pada Listing~\ref{lst:struktur-gfd}, baru \texttt{config.py} dan \texttt{dataset/} yang ada berkasnya.
Lima subpaket sisanya masih kosong.

\section{Akuisisi per Sumber}
\label{sec:akuisisi}

\subsection{PLN Puslitbang}
\label{subsec:pln}
Data sambaran PLN Puslitbang tidak memiliki jalur pengunduhan.
Rekaman itu diperoleh melalui pembimbing dalam bentuk berkas, dan datanya tidak terbuka untuk umum, sehingga tidak dapat disebarluaskan bersama laporan ini.
Kode hanya memuat pembacaan berkasnya, tanpa bagian akuisisi.
Berkas yang diterima berbentuk lembar kerja, satu baris satu sambaran.
Isi kolomnya dapat dilihat pada Tabel~\ref{tab:kolom-pln}.

\begin{table}[!t]
  \caption{Kolom berkas sambaran PLN Puslitbang}
  \label{tab:kolom-pln}
  \footnotesize
  \begin{tabular}{ | l | p{0.55\textwidth} | }
    \hline
    Kolom & Keterangan \\
    \hline
    Select & Penanda pilihan bawaan antarmuka ekspor; tidak dipakai \\
    Date and time & Tanggal dan waktu kejadian sambaran, sampai milidetik \\
    Latitude & Koordinat lintang sambaran, dalam derajat \\
    Longitude & Koordinat bujur sambaran, dalam derajat \\
    Discrimination & Jenis sambaran, misalnya CG$-$ untuk sambaran awan ke tanah bermuatan negatif \\
    Signal (kA) & Arus puncak sambaran dalam kiloampere; tandanya menunjukkan polaritas \\
    Chi square & \textcolor{red}{[PERLU KONFIRMASI: ukuran kecocokan penyelesaian lokasi]} \\
    Ell. smaj & Sumbu panjang elips galat lokasi \\
    Ell. smin & Sumbu pendek elips galat lokasi \\
    Ell. angle & Sudut orientasi elips galat lokasi, dalam derajat \\
    Freedom & \textcolor{red}{[PERLU KONFIRMASI: derajat kebebasan penyelesaian lokasi]} \\
    Multi. & \textcolor{red}{[PERLU KONFIRMASI: jumlah sambaran dalam satu kilat]} \\
    Sensors nb. & Jumlah sensor yang mendeteksi sambaran \\
    \hline
  \end{tabular}
\end{table}

\subsection{MERLIN}
\label{subsec:merlin}
Data sambaran MERLIN diambil dari arsip cuaca milik Kennedy Space Center \autocite{ksc2026merlin}, yang dibahas pada Subbab~\ref{subsec:sumber-merlin}.
Arsip itu tidak menyediakan API yang terdokumentasi, sehingga jalur pengunduhan disusun dengan meniru permintaan yang dipakai antarmuka webnya.
Permintaan tersebut memerlukan token terkode.
Penyandian token itu disusun ulang di dalam kode, lalu hasilnya dicocokkan dengan tujuh token yang ditangkap dari antarmuka aslinya, dan ketujuhnya cocok.

Rentang waktu satu permintaan dibatasi oleh arsip, sehingga periode 2018 sampai 2024 dipecah menjadi beberapa jendela unduhan.
Arsip MERLIN juga dibatasi secara geografis, dan permintaan dari luar wilayah yang diizinkan akan ditolak.
Pembatasan ini diatasi dengan VPN yang bekerja pada tingkat perangkat, bukan dengan ekstensi peramban, karena permintaan dikirim oleh kode dan bukan oleh peramban.

Berkas hasil unduhan berbentuk CSV, satu baris satu sambaran.
Isi kolomnya dapat dilihat pada Tabel~\ref{tab:kolom-merlin}.

\begin{table}[!t]
  \caption{Kolom berkas sambaran MERLIN}
  \label{tab:kolom-merlin}
  \footnotesize
  \begin{tabular}{ | l | p{0.55\textwidth} | }
    \hline
    Kolom & Keterangan \\
    \hline
    Date & Tanggal kejadian sambaran \\
    Time & Waktu kejadian sambaran, sampai pecahan detik \\
    Latitude & Koordinat lintang sambaran, dalam derajat \\
    Longitude & Koordinat bujur sambaran, dalam derajat \\
    Signal Strength & Kuat sinyal sambaran; tandanya menunjukkan polaritas \\
    SemiMajor Axis 50\% CI & Sumbu panjang elips galat lokasi pada selang kepercayaan 50\% \\
    SemiMinor Axis 50\% CI & Sumbu pendek elips galat lokasi pada selang kepercayaan 50\% \\
    Ellipse Angle & Sudut orientasi elips galat lokasi, dalam derajat \\
    Sensors & Daftar nomor sensor yang mendeteksi sambaran \\
    \hline
  \end{tabular}
\end{table}

\subsection{ERA5}
\label{subsec:era5}
Data ERA5 diambil dari \textit{Climate Data Store} milik Copernicus melalui pustaka \texttt{cdsapi} \autocite{c3s2023era5}.
ERA5 adalah reanalisis, yaitu keluaran model cuaca yang digabungkan dengan hasil pengamatan di seluruh dunia melalui asimilasi data, sehingga menghasilkan medan yang lengkap dan konsisten.
Produk yang dipakai adalah ERA5 \textit{single levels} pada resolusi per jam, yang disajikan pada grid lintang-bujur beraturan berukuran 0,25\textdegree.
Variabel yang diminta dipilih berdasarkan parameter meteorologis yang dibahas pada Subbab~\ref{subsec:parameter-meteorologis}.
Pemetaan tiap parameter ke variabelnya dapat dilihat pada Tabel~\ref{tab:era5-variabel}.

Permintaan dibatasi pada kotak koordinat tiap domain, bukan seluruh dunia.
Kotak itu diturunkan dari sebaran koordinat data sambaran PLN dan MERLIN, lalu dilebarkan ke luar sampai batas sel grid terdekat.

\begin{table}[!t]
  \caption{Pemetaan parameter meteorologis ke variabel ERA5}
  \label{tab:era5-variabel}
  \footnotesize
  \begin{tabular}{ | l | p{0.26\textwidth} | p{0.40\textwidth} | }
    \hline
    Kolom & Parameter & Variabel ERA5 \\
    \hline
    CAPE & Energi potensial konvektif & convective\_available\_\newline potential\_energy \\
    KX & Indeks K & k\_index \\
    TCIW & Kandungan es awan satu kolom & total\_column\_cloud\_ice\_water \\
    TCLW & Kandungan air cair awan satu kolom & total\_column\_cloud\_liquid\_\newline water \\
    VIIWD & Divergensi fluks air beku awan & vertical\_integral\_of\_divergence\_\newline of\_cloud\_frozen\_water\_flux \\
    VILWD & Divergensi fluks air cair awan & vertical\_integral\_of\_divergence\_\newline of\_cloud\_liquid\_water\_flux \\
    VIMDF & Divergensi fluks kelembapan & vertical\_integral\_of\_divergence\_\newline of\_moisture\_flux \\
    TOTALX & Indeks \textit{total totals} & total\_totals\_index \\
    CIN & Penghambat konveksi & convective\_inhibition \\
    CBH & Tinggi dasar awan & cloud\_base\_height \\
    TCWV & Uap air satu kolom & total\_column\_water\_vapour \\
    D2M & Suhu titik embun pada 2 meter & 2m\_dewpoint\_temperature \\
    CRR & Laju hujan konvektif & convective\_rain\_rate \\
    \hline
  \end{tabular}
\end{table}

Nilai ERA5 yang dipakai bersifat sesaat, bukan rerata.
Setiap medan berlaku tepat pada awal jam, bukan sebagai rata-rata sepanjang jam tersebut.
Karena target menghitung sambaran di dalam rentang satu jam, pemasangan prediktor dengan target mengikuti konvensi pada Subbab~\ref{subsec:grid-temporal}.

Seluruh unduhan mencakup setiap jam sepanjang 2018 sampai 2024 untuk kedua domain, yaitu 61.368 jam per domain.
Pengunduhan dilakukan dalam potongan bulanan, satu permintaan untuk satu domain satu bulan, karena permintaan satu domain untuk satu tahun penuh melewati batas biaya permintaan CDS.

\subsection{NASA POWER}
\label{subsec:power}
Data NASA POWER diambil dari layanan API yang disediakan proyek POWER melalui permintaan HTTP \autocite{nasapower}.
POWER adalah proyek NASA yang menyediakan data meteorologis dan surya dalam bentuk yang siap dipakai ulang.
Parameter meteorologisnya berasal dari MERRA-2, yaitu reanalisis milik NASA, dan disajikan pada grid berukuran 0,5\textdegree{} lintang kali 0,625\textdegree{} bujur.
Parameter yang diminta dipilih berdasarkan pembahasan pada Subbab~\ref{subsec:parameter-meteorologis}, dan pemetaannya dapat dilihat pada Tabel~\ref{tab:power-variabel}.

Titik pengambilan diturunkan dari kotak koordinat yang sama dengan ERA5, tetapi bentuk permintaannya berbeda.
Grid proyek berukuran 0,5\textdegree{} pada kedua arah, seperti dijelaskan pada Subbab~\ref{subsec:grid-spasial}, sedangkan grid POWER berukuran 0,5\textdegree{} lintang dan 0,625\textdegree{} bujur.
Langkah bujur POWER lebih besar, sehingga terdapat sel proyek yang berbeda tetapi titik POWER terdekatnya sama.

Layanan POWER pada resolusi per jam hanya melayani permintaan untuk satu titik, bukan untuk satu wilayah, sehingga titik yang diminta harus ditentukan lebih dahulu.
Yang diminta adalah titik grid asli POWER, bukan titik tengah tiap sel proyek.
Jika permintaan disusun per sel proyek, sel yang berbagi titik POWER yang sama akan menghasilkan dua permintaan untuk data yang sama.

\begin{table}[!t]
  \caption{Pemetaan parameter meteorologis ke variabel NASA POWER}
  \label{tab:power-variabel}
  \footnotesize
  \begin{tabular}{ | l | p{0.26\textwidth} | p{0.40\textwidth} | }
    \hline
    Kolom & Parameter & Variabel NASA POWER \\
    \hline
    PRECTOTCORR & Presipitasi terkoreksi & PRECTOTCORR \\
    PS & Tekanan permukaan & PS \\
    T2M & Suhu udara pada 2 meter & T2M \\
    RH2M & Kelembapan relatif pada 2 meter & RH2M \\
    WS2M & Kecepatan angin pada 2 meter & WS2M \\
    \hline
  \end{tabular}
\end{table}

Satu parameter POWER yang dipakai penelitian pendahulu tidak diminta pada tugas akhir ini, yaitu AOD\_55\_ADJ.
POWER menerbitkan parameter itu pada resolusi bulanan dan tidak lebih halus.
Pada resolusi per jam nilainya tetap sepanjang satu bulan penuh, sehingga tidak dapat menjelaskan ragam antarjam pada target.
Alasan penggugurannya dibahas pada Subbab~\ref{subsec:parameter-meteorologis}.

\section{Pemuatan, Penyelarasan, dan Penggabungan}
\label{sec:penggabungan}
Keempat sumber memakai satuan ruang dan waktu yang berbeda, sehingga tidak dapat langsung digabungkan.
Subbab ini menjelaskan penyelarasan yang dilakukan lebih dahulu, lalu cara keempatnya disatukan menjadi satu tabel.

Penyelarasan dilakukan terhadap satu kunci yang sama, yaitu lintang sel, bujur sel, dan jam.
Rekaman sambaran dihitung jumlahnya per sel per jam.
Nilai ERA5 dirata-ratakan dari titik aslinya yang berukuran 0,25\textdegree{} ke sel proyek yang berukuran 0,5\textdegree.
Nilai POWER dipetakan dari titik gridnya ke setiap sel yang memakai titik tersebut.
Waktu dari semua sumber diubah ke UTC lebih dahulu, kemudian penandanya dilepas, sehingga seluruh sumber memakai kalender yang sama.

Rangka tabel dibentuk lebih dahulu, dan pada tahap itu belum memuat data apa pun.
Rangka tersebut adalah hasil perkalian setiap sel grid dengan setiap jam pada rentang yang dikonfigurasi.
Jumlah sambaran per sel-jam kemudian ditempelkan ke rangka itu dengan \textit{left join}, dan sel-jam yang tidak memiliki sambaran diisi nol.
Dengan cara ini, sel-jam tanpa sambaran tetap ada sebagai baris yang lengkap, bukan menghilang dari tabel.
Nilai POWER dan ERA5 ditempelkan ke tabel yang sama dengan cara yang sama.

Penggabungan semacam ini punya satu kegagalan yang tidak bersuara.
Penggabungan yang tidak mencocokkan satu kunci pun tidak memunculkan galat, melainkan menghasilkan tabel utuh yang seluruh prediktornya kosong.
Karena itu pemeriksaan kunci dijalankan sebelum tiap penggabungan, dan pemeriksaan berhenti dengan galat bila tipe data kunci berbeda atau irisan kunci sama sekali kosong.

Terdapat tiga hal yang dapat membuat irisan kunci kosong, dan tidak satu pun memunculkan galat dengan sendirinya.
Pertama, tipe data pada kunci waktu berbeda antara sumber dan rangka.
Kedua, sumber memakai acuan jam yang berbeda, sehingga jamnya bergeser.
Ketiga, bujur ERA5 datang pada rentang 0 sampai 360 derajat, sedangkan data sambaran memakai rentang $-$180 sampai 180 derajat.
Sebab ketiga hanya menimpa domain subtropis, karena bujur Jawa Barat bernilai positif pada kedua rentang tersebut.